\documentclass[../../main.tex]{subfiles}
\graphicspath{{../../image/}{../../figure/}}
\begin{document}

\section{弱前沿的爆发}\label{section:ch07-弱前沿爆发}

本节旨在揭示真正非线性性在弱前沿放大并最终爆发中所起的决定性作用。

我们考虑严格双曲方程组\eqref{align:ch07-典范守恒律}的一个 Lipschitz 连续解 $U$，定义在带形 $(-\infty,\infty)\times[0,T)$ 上，并具有如下结构：一条 $C^1$ 曲线 $x=\chi(t)$ 从原点发出，在集合 $\{(x,t):0\leqslant t<T,\ x>\chi(t)\}$ 上 $U(x,t)=\bar{U}=\text{常数}$；而在集合 $\{(x,t):0\leqslant t<T,\ x<\chi(t)\}$ 上 $U$ 是 $C^2$，其一阶与二阶偏导数当 $x\uparrow\chi(t)$ 时达到非零极限。因而 $\chi(\cdot)$ 是传入常状态的弱前沿。

在集合 $\{(x,t):0\leqslant t<T,\ x<\chi(t)\}$ 上，
\begin{align}
\partial_t U(x,t)+\mathrm{D}F(U(x,t))\partial_x U(x,t)=0.
\label{align:ch07-弱前沿方程}
\end{align}
由于 $U(\chi(t)-,t)=\bar{U}$，
\begin{align}
\partial_t U(\chi(t)-,t)+\dot{\chi}(t)\partial_x U(\chi(t)-,t)=0.
\label{align:ch07-弱前沿跃动}
\end{align}
把\eqref{align:ch07-弱前沿方程}与\eqref{align:ch07-弱前沿跃动}结合，
\begin{align}
[\mathrm{D}F(\bar{U})-\dot{\chi}(t)I]\partial_x U(\chi(t)-,t)=0.
\label{align:ch07-弱前沿特征值}
\end{align}
因而 $\dot{\chi}(t)$ 是常数，等于某特征族 $i$ 的 $\lambda_i(\bar{U})$，且
\begin{align}
\partial_x U(\chi(t)-,t)=a(t)R_i(\bar{U}).
\label{align:ch07-弱前沿振幅}
\end{align}
函数 $a(t)$ 度量弱前沿的强度。

我们从左用 $L_i(\bar{U})$ 乘\eqref{align:ch07-弱前沿振幅}，利用\eqref{align:ch07-左右向量规范化}并关于 $t$ 求导，得到
\begin{align}
\frac{\mathrm{d}a(t)}{\mathrm{d}t}=L_i(\bar{U})\left[\partial_x\partial_t U(\chi(t)-,t)+\lambda_i(\bar{U})\partial_x\partial_x U(\chi(t)-,t)\right].
\label{align:ch07-da1}
\end{align}

接下来，从左用 $L_i(U(x,t))$ 乘\eqref{align:ch07-弱前沿方程}，
\begin{align}
L_i(U(x,t))\left[\partial_t U(x,t)+\lambda_i(U(x,t))\partial_x U(x,t)\right]=0,
\label{align:ch07-弱前沿左乘}
\end{align}
然后关于 $x$ 求导并令 $x\uparrow\chi(t)$。合并 $\dot{\chi}(t)=\lambda_i(\bar{U})$、\eqref{align:ch07-弱前沿跃动}、\eqref{align:ch07-da1}、\eqref{align:ch07-弱前沿振幅}与\eqref{align:ch07-左右向量规范化}，我们结论 $a(t)$ 满足一个 Bernoulli 型常微分方程：
\begin{align}
\frac{\mathrm{d}a}{\mathrm{d}t}+\mathrm{D}\lambda_i(\bar{U})R_i(\bar{U})a^2=0.
\label{align:ch07-Bernoulli}
\end{align}

因而，若 $\bar{U}$ 是 $i$-特征族的真正非线性态且 $\mathrm{D}\lambda_i(\bar{U})R_i(\bar{U})a(0)<0$，则弱波的强度随时间增长，并最终当 $t\uparrow[-\mathrm{D}\lambda_i(\bar{U})R_i(\bar{U})a(0)]^{-1}$ 时爆发。古典解破裂的问题将在下一节从更广阔的视角加以讨论。

\end{document}
